\documentclass[11pt]{article}

\usepackage[margin=1in]{geometry}
\usepackage{amsmath,amssymb,amsthm,mathtools}
\usepackage{enumitem}
\usepackage{booktabs}
\usepackage{tabularx}
\usepackage{microtype}
\usepackage[hidelinks]{hyperref}
\usepackage{xurl}

\newtheorem{definition}{Definition}
\newtheorem{remark}{Remark}

\newcommand{\MaxL}{\mathcal{L}}

\title{Receipt Line Items for Anonymity Claims:\\
\large A Minimal Tuple Schema for Published Budgets (Series Synthesis)}
\author{}
\date{Unified archive note v0.31 (2026-03-05; archive v162)}

\begin{document}
\maketitle

\begin{abstract}
Many papers in this archive culminate in a short published object: a label (OINL), a budget (e.g., MaxL), and a receipt that pins the claim to concrete knobs under change control.
(Receipt envelopes are specified by ABOM/OINL in Anonymity~C; this note standardizes the line-item tuples carried inside.)\cite{series:abomoinl}

To keep the series non-redundant and referee-grade, we centralize the receipt-facing interface:
\emph{what minimal tuple must appear as a line item so that a reader can understand the claim and an auditor can request exactly the missing artifacts?}
We give a small schema that (i) is stable across modules (retries, congestion, destination equalization, state/auditing), and (ii) cleanly separates \emph{budget semantics} from \emph{publication mechanics}.
\end{abstract}

\paragraph{Scope.}
This is a short schema/interface note.
It does not propose new mechanisms; it standardizes \emph{what to publish} when a mechanism paper exports a budget.

\paragraph{Imports.}
Contract-object vocabulary is imported from Synthesis~0.\cite{series:contractobjects}
ABOM/OINL packaging is imported from Anonymity~C.\cite{series:abomoinl}
Transparency primitives (beacon/VRF binding, append-only logs, inclusion/consistency proofs) are imported from Synthesis~9.\cite{series:transparencyprimitives}
Endpoint metric semantics (odds inflation / MaxL and conversions) are imported from Synthesis~5 and Anonymity~A.\cite{series:endpointbridge,series:oddsinflation}
Replay-plan and state-declaration drift/equivalence rules are imported from Synthesis~18.\cite{series:planstateequiv}

\paragraph{Exports.}
A minimal, receipt-ready \emph{line item tuple} and a ``who owns what'' map so other papers can remain tight.

\section{Receipt envelope (three IDs + optional pointers)}
A receipt is a small published object intended to be stable across modules.
To keep replay and drift semantics non-redundant, we separate the \emph{envelope} (series-wide identifiers) from per-claim line items.

\begin{definition}[Receipt envelope]
A receipt begins with:
\begin{itemize}[leftmargin=*]
\item \texttt{claim\_id}: a stable identifier for the claim family (often a digest ID);
\item \texttt{tw\_id}: the threat-window declaration governing repetitions/resets;
\item \texttt{state\_decl\_id} (optional): scoped state/measurement declaration used during evaluation; include when long-lived state or monitored quantities enter the guarantee.
\end{itemize}
Identifier and digest conventions are centralized in Synthesis~8; plan/state equivalence is centralized in Synthesis~18.\cite{series:notation,series:planstateequiv}
\end{definition}

\begin{remark}[Why the envelope matters]
Most ``anonymity bugs'' in deployed systems are drift bugs: the receipt is still published, but the knobs/state/measurement surface has changed.
Separating the envelope makes it explicit what must be re-issued (new \texttt{state\_decl\_id} or \texttt{tw\_id}) versus what may be replayed (same \texttt{plan\_spec\_id} under Synthesis~18).
\end{remark}

\paragraph{Auditor replay.}
This note defines the tuple surface; the canonical mechanical replay checklist (anchoring $\rightarrow$ bundle fetch $\rightarrow$ to-spec replay) is Synthesis~20.\cite{series:auditorreplay}

\section{The line item (definition)}
\begin{definition}[Receipt line item]
A \emph{receipt line item} is a short, hash-addressable record that binds a \emph{budget claim} to:
(i) a declared witness surface (transcript tier/projection),
(ii) a declared repetition window and usage assumption, and
(iii) the specific knobs/parameters whose drift would change the claim.
The item is designed so an auditor can replay the accountant given access to the artifact bundle referenced by ABOM/OINL.
\end{definition}


\paragraph{Contact surfaces are semantic.}
When a line item names an external service contact surface (e.g., delegated routing or a fallback committee), the contact surface is a canonical description of the interface: \emph{method}, full URL (including path scope), and any declared request-shaping defaults (e.g., IPIP-0484 protocol filtering) and time budgets that affect who is contacted and how termination-time behaves. Changing host, path, method, default filter parameters, or declared time budgets is interface drift and triggers recertification.
\cite{series:deployedsurfaces,series:planstateequiv,series:recert}

\begin{remark}
For publication and user micro-diff hygiene, we recommend publishing \texttt{contact\_surface\_id} as \texttt{sha256(JCS(contact\_surface))} (JCS = RFC~8785)\cite{rfc8785} so the interface can be compared atomically across releases.
\end{remark}

\section{Minimal tuple schema (stable across modules)}
We recommend the following minimal fields.
(Names are schematic; deployments may choose concrete encodings, but the semantics should match.)

\begin{table}[t]
\centering
\begin{tabularx}{\linewidth}{@{}lX@{}}
\toprule
Field & Meaning (and who owns the semantics) \\
\midrule
\texttt{claim\_id} & Stable identifier for the claim object (ABOM/OINL carrier).\cite{series:abomoinl} \\
\texttt{exposure\_nf\_id} & Stable exposure normal-form id / schema hash for the projection being budgeted. This is the real interface identifier; the rest of the tuple refines it.\cite{series:traceifaces} \\
\texttt{witness\_tier} (optional) & Human-readable tier label for quick reading; redundant to \texttt{exposure\_nf\_id} and should be treated as non-binding.\cite{series:traceifaces} \\
\texttt{tw\_id} & Threat-window declaration: what repetition horizon the claim is sized for.\cite{series:threatwindows} \\
\texttt{state\_decl\_id} (optional) & Digest-bound state-surface declaration committing to which long-lived control-plane state surfaces are in scope for the bound (and which are excluded). This includes deployed router sets, path scope, request-shaping defaults (e.g., filters), timeout/refresh cadences, downgrade/fallback rules, and retrieval-mode switches when they affect the transcript.\cite{series:state1,series:deployedsurfaces,series:workedexample} \\
\texttt{contact\_surface\_id} (optional) & Digest-bound id for a named external service contact surface (method + URL + request-shaping defaults + time budgets). Recommended: \texttt{sha256(JCS(contact\_surface))}.\cite{rfc8785} \\
\texttt{evidence\_id} (optional) & Digest-bound identifier for a concrete evidence object (log audit table, dual certificate, proof transcript, \dots) supporting \texttt{budget\_value}. Naming conventions live in Synthesis~8; conditional/per-step evidence objects are standardized in Synthesis~21.\cite{series:notation,series:condcert} \\

\texttt{usage} & Usage assumption inside the window (e.g., $Q$ lookups/window; rate-limit policy).\cite{series:bossfightB} \\
\texttt{budget\_type} & What is being bounded (e.g., MaxL in bits, OINL odds inflation, or a TV certificate).\cite{series:endpointbridge,series:knapsacktransport} \\
\texttt{budget\_value} & The bound value (e.g., $b$ bits per lookup, or an identified OINL label).\cite{series:oddsinflation,series:abomoinl} \\
\texttt{obs\_model} & Observation parameterization (e.g., $p_{\mathrm{obs}}$, a monitored upper bound, a tiered observation vector, or tiered path-selection bounds for fallback use).\cite{series:obsatten,series:tieredobs,series:fallbacktax} \\
\texttt{knobs} & Hashes/values of the knobs that determine the accountant (retries, timeouts, widths, calendars, equalization params).\cite{series:bossfightB,series:mceq} \\
\texttt{lineage} & Optional predecessor pointer plus relation tag (\texttt{same-enf}, \texttt{coarsening}, or \texttt{new-interface}). This is what makes cross-release reuse explicit instead of implicit.\cite{series:traceifaces,series:recert,series:releaseproperty} \\
\texttt{replay\_hook} & How an auditor replays the accountant: a stable \texttt{plan\_id} plus a digest-bound \texttt{plan\_spec\_id} (canonical plan specification hash) and an artifact-bundle pointer (on request).\cite{series:contractobjects,series:transparencyprimitives} \\
\bottomrule
\end{tabularx}
\caption{A minimal receipt line-item schema. Papers should avoid restating semantics for these fields; instead, cite the owning note listed in the right column.}
\label{tab:minschema}
\end{table}
\begin{remark}[Evidence \emph{versus} replay]
\label{rem:evidence-vs-replay}
\texttt{evidence\_id} names \emph{what} supports the number in \texttt{budget\_value} (a log-audit table, a dual certificate, a proof transcript, \dots).
\texttt{replay\_hook} names \emph{how} an auditor should recompute and validate it (a stable \texttt{plan\_id} with a digest-bound \texttt{plan\_spec\_id}, plus the associated bundle pointer on request).
They are intentionally separate: a single replay plan may emit multiple evidence objects, and an evidence object may be checkable by more than one replay plan.
\end{remark}




\begin{remark}[Separation-class primary paths need an explicit guard]
Some deployments publish a primary-path line item whose numeric value is interpreted only under a role-separation assumption (e.g., ``0 bits on a linkability surface under \texttt{sep.nr1}'').
In that case, the line item should commit to a separation-manifest digest (Synthesis~22) inside \texttt{obs\_model} or \texttt{knobs} so the non-collusion assumption and invalidators are diffable.
If the deployment is willing to publish an upper bound $\rho$ on separation failure within the threat window, carry it as \texttt{rho\_upper} in the same place and interpret it as the guarded checkpoint $(0,\rho)$ per Synthesis~23.
Tightening \texttt{rho\_upper} should be treated as replay/evidence-required; relaxing it may be tagged as \texttt{lineage=coarsening} but still must be published and diffed.\cite{series:primarysepmanifest,series:probsep}
\end{remark}

\begin{remark}[Why this is the right granularity]
The tuple intentionally stops short of shipping full artifacts: it makes the claim \emph{addressable}.
A reader can compare releases from the public tuple alone, while an auditor can request the precise evidence referenced by \texttt{replay\_hook}.
This matches the source-only artifact policy.\cite{series:artifactpolicy}
If deployments want end users to \emph{discover} the right receipt and mechanically validate its anchoring, pair the tuple with the minimal user-verifiability record in Synthesis~13.\cite{series:userminiface}
\end{remark}

\begin{remark}[Interface ids and lineage prevent accidental claim reuse]
A line item should never rely on prose like ``same mechanism as last release.''
The pair $(\texttt{exposure\_nf\_id},\texttt{lineage})$ makes the intended relation explicit:
if the interface is unchanged, the new receipt says so; if the interface is an explicit coarsening, the receipt says so; otherwise the receipt declares a new interface and stops pretending the old certificate applies automatically.
This keeps release engineering, recertification, and user-facing validation aligned.

Replay workflows and state assumptions have the same ``label drift'' failure mode: a bare plan name or a human state label is not a stable audit promise.
Receipts should therefore carry digest-bound \texttt{plan\_spec\_id} and (when relevant) \texttt{state\_decl\_id}, and treat changes to these ids using the equivalence/drift taxonomy centralized in Synthesis~18.\cite{series:planstateequiv}
\end{remark}

\section{Module-specific instantiations (short)}
This note intentionally avoids module-specific prose.
If a mechanism needs extra fields beyond Table~\ref{tab:minschema}, publish them under a single \texttt{module\_ext} map and bind any semantic interfaces (e.g., service contact surfaces) by digest ids such as \texttt{contact\_surface\_id}.
A module that exports a receipt line item should only specify what it \emph{owns} (the meaning of \texttt{usage}, \texttt{obs\_model}, \texttt{knobs}, \texttt{budget\_type}, and any evidence objects) and cite this note for the tuple shape.
For an end-to-end instantiation (including the interplay with state declarations and deployed routing knobs), see the worked example (Synthesis~17) and the deployed-surface note (Synthesis~30).\cite{series:workedexample,series:deployedsurfaces}\section{A publication checklist (one paragraph)}
To publish a line item responsibly:
(i) include Table~\ref{tab:minschema}'s fields;
(ii) include a stable \texttt{exposure\_nf\_id} and, when applicable, a \texttt{lineage} relation to the previous claim;
(iii) commit the tuple to ABOM/OINL and to a release receipt;\cite{series:abomoinl,series:releaseproperty}
(iv) anchor it in a transparency log with inclusion/consistency proofs;\cite{series:transparencyprimitives}
and (v) treat any knob drift or interface drift as a recertification trigger (Certified~C).\cite{series:recert}
When public discovery matters, publish a companion UVI record that points users to the anchored receipt tuple and its validation path.\cite{series:userminiface}
For a concrete end-to-end worked instantiation of these fields, see the worked-example micro-note (Synthesis~17).\cite{series:workedexample}

\begin{thebibliography}{99}\small


\bibitem{rfc8785}
A.~Rundgren.
\newblock JSON Canonicalization Scheme (JCS).
\newblock RFC~8785, IETF, 2020.

\bibitem{series:contractobjects}
Anonymous.
\newblock Contract Objects for Anonymous-DHT Privacy Budgets and Claims.
\newblock Synthesis~0, The-One-True-Archive (2026).

\bibitem{series:abomoinl}
Anonymous.
\newblock ABOM and OINL for Anonymity Claims.
\newblock Anonymity series Paper~C, The-One-True-Archive (2026).

\bibitem{series:transparencyprimitives}
Anonymous.
\newblock Transparency Primitives for Privacy Receipts: Beacons, VRFs, and Append-Only Logs.
\newblock Synthesis~9, The-One-True-Archive (2026).

\bibitem{series:endpointbridge}
Anonymous.
\newblock Endpoint Metrics Bridge: Odds Inflation, PML/MaxL, and Identification Sizing.
\newblock Synthesis~5, The-One-True-Archive (2026).

\bibitem{series:oddsinflation}
Anonymous.
\newblock Odds Inflation, Privacy Loss, and Endpoint Anonymity Bounds.
\newblock Anonymity series Paper~A, The-One-True-Archive (2026).

\bibitem{series:traceifaces}
Anonymous.
\newblock Trace Interfaces for Anonymous DHTs.
\newblock Synthesis~3, The-One-True-Archive (2026).

\bibitem{series:threatwindows}
Anonymous.
\newblock Threat Models and Repetition Windows for Anonymous DHTs.
\newblock Synthesis~4, The-One-True-Archive (2026).

\bibitem{series:planstateequiv}
Anonymous.
\newblock Digest-Bound Replay Plans and State Declarations: Equivalence, Minimality, and Change Control.
\newblock Synthesis~18, The-One-True-Archive (2026).

\bibitem{series:auditorreplay}
Anonymous.
\newblock Auditor Replay for Anonymous-DHT Receipt Line Items: A Mechanical Checklist for Digest-Bound Claims.
\newblock Series synthesis note (Synthesis~20), 2026. (This archive.)

\bibitem{series:fallbacktax}
Anonymous.
\newblock Adaptive Fallback as a Witness Stream: The Path-Selection Tax and Honest Tiered Receipts.
\newblock Series synthesis note (Synthesis~14), The-One-True-Archive (2026).

\bibitem{series:obsatten}
Anonymous.
\newblock Observation Attenuation for Repeated-Use Privacy Budgets: Mixture/Subsampling Bounds and a Stub-Mix Interface.
\newblock Synthesis~11, The-One-True-Archive (2026).

\bibitem{series:tieredobs}
Anonymous.
\newblock Tiered Observation Vectors for Anonymous-DHT Receipts: Contact-Set, Timing, and Congestion Tiers with Conservative MaxL Composition.
\newblock Synthesis~15, The-One-True-Archive (2026).

\bibitem{series:bossfightB}
Anonymous.
\newblock AnonDHT in the Boss Fight Regime: Compiling a MaxL Budget into Lookup Knobs.
\newblock Boss Fight series Paper~B, 2026. (This archive.)

\bibitem{series:mceq}
Anonymous.
\newblock MC-EQ and CoverSketch: Destination Privacy via Markov-Chain Equalization in Anonymous DHT Lookups.
\newblock Release \& destination Paper~B, The-One-True-Archive (2026).

\bibitem{series:wcongeq}
Anonymous.
\newblock W-Congestion-EQ: Verifiable Congestion Budgets (Certificates + Receipts).
\newblock Congestion series Paper~3, The-One-True-Archive (2026).

\bibitem{series:knapsacktransport}
Anonymous.
\newblock Optimal Poset-Feasible Padding: Knapsack, Transport, and Dual Certificates for TV Privacy.
\newblock Published backbone note (2026-01-29). Wiki: \texttt{[[2026.01.29 - Mathematics: Optimal Poset-Feasible Padding: Knapsack, Transport, and Dual Certificates for TV Privacy]]}.

\bibitem{series:state3}
Anonymous.
\newblock Closed-View Auditing and CPPC Manifests.
\newblock AnonDHT state series Paper~3, The-One-True-Archive (2026).

\bibitem{series:userminiface}
Anonymous.
\newblock User-Verifiable Receipt Interfaces for Anonymous-DHT Claims: A Minimal Discoverability and Validation Contract.
\newblock Synthesis~13, The-One-True-Archive (2026).

\bibitem{series:primarysepmanifest}
Anonymous.
\newblock Primary-Path Separation Manifests: Guard Objects for Non-Collusion Assumptions in Reader-Privacy Receipts.
\newblock Synthesis~22, The-One-True-Archive (2026).

\bibitem{series:probsep}
Anonymous.
\newblock Probabilistic Separation Budgets: Separation Failure as a Receipt-Grade $(\varepsilon,\delta)$ Adjunct.
\newblock Synthesis~23, The-One-True-Archive (2026).

\bibitem{series:condcert}
Anonymous.
\newblock Conditional Per-Step Budget Certificates: A Receipt-Grade Interface for Adaptive Horizon Claims.
\newblock Synthesis~21, The-One-True-Archive (2026).

\bibitem{series:workedexample}
Anonymous.
\newblock Worked Example: Instantiating a Tiered Reader-Privacy Receipt.
\newblock Synthesis~17, The-One-True-Archive (2026).

\bibitem{series:artifactpolicy}
Anonymous.
\newblock Artifact and Disclosure Policy for the One-True-Archive (Source-Only Public Release).
\newblock Synthesis~6, The-One-True-Archive (2026).

\bibitem{series:releaseproperty}
Anonymous.
\newblock Privacy as a Release Property: Proof-Carrying Anonymity Receipts for Anonymous DHTs.
\newblock Release \& destination Paper~A, The-One-True-Archive (2026).

\bibitem{series:recert}
Anonymous.
\newblock Disagreement-Gated Recertification.
\newblock Certified series Paper~C, The-One-True-Archive (2026).


\bibitem{series:notation}
Anonymous.
\newblock Notation and Minimal Model for the One-True-Archive.
\newblock Synthesis~8, The-One-True-Archive (2026).


\bibitem{series:state1}
Anonymous.
\newblock State-Dependent Anonymity: Selection and State Evolution as an Attack Surface.
\newblock AnonDHT State series Paper~1, The-One-True-Archive (2026).


\bibitem{series:deployedsurfaces}
Anonymous.
\newblock Deployed Observation Surfaces for Anonymous DHTs: Control-plane knobs in IPFS/libp2p delegated routing and OHTTP evolution.
\newblock Series synthesis note (Synthesis~30), 2026. (This archive.)

\end{thebibliography}

\end{document}